\section{问题二：基于离散制氨调节的运行优化}

\subsection{问题分析与求解思路}

问题二将园区制氨产能增至72吨/日，电氢氨装置仅有全额开机和停机两种运行方式。需要确定不同产量水平（72、63、54、45、36吨/日）下成本最低的生产时段安排。产能扩容后设备参数线性提升：ALKEL 20 MW、PEMEL 20 MW、合成氨装置1.5 MW，电氢氨总额定功率$P_{H2NH3}^{max} = 41.5$ MW。每种产量对应的开机时数为$k = Q_{NH3} / 3.0$小时。

\begin{figure}[H]
\centering
\includegraphics[width=0.85\textwidth,height=0.38\textheight,keepaspectratio]{../figures/fig_flow_q2.pdf}
\caption{问题二求解流程图}\label{fig:flow_q2}
\end{figure}

本问题的核心是一个0-1整数规划（MILP）问题：在24个时段中选择$k$个时段开机运行，使得日运行总成本最小（图\ref{fig:flow_q2}）。由于决策变量为二进制变量，约束和目标函数均为线性，可采用分支定界法精确求解\cite{ref2}。此类生产调度问题在电力系统机组组合领域已有成熟的MILP建模方法\cite{ref3}。

\subsection{数学模型建立}

\subsubsection{决策变量}

引入二进制决策变量$u(t) \in \{0, 1\}$，$t = 1, 2, \ldots, 24$，其中$u(t) = 1$表示时段$t$电氢氨装置全额开机运行，$u(t) = 0$表示停机。

\subsubsection{约束条件}

产量约束要求总开机时数等于目标产量对应的小时数：
\begin{equation}\label{eq:q2_production}
\sum_{t=1}^{24} u(t) = k = \frac{Q_{NH3}}{r_{NH3}} = \frac{Q_{NH3}}{3.0}
\end{equation}

功率平衡约束为：
\begin{equation}
P_{buy}(t) - P_{sell}(t) = P_{load}(t) + u(t) \cdot P_{H2NH3}^{max} - P_{wind}(t) - P_{pv}(t)
\end{equation}

购售电非负约束：$P_{buy}(t) \geq 0$，$P_{sell}(t) \geq 0$。

\subsubsection{目标函数}

最小化日运行总成本：
\begin{equation}\label{eq:q2_obj}
\min \quad C_{day} = C_{RE} + \sum_{t=1}^{24} \lambda_{buy}(t) \cdot P_{buy}(t) \cdot \Delta t + \sum_{t=1}^{24} u(t) \cdot C_{OM,h} - \lambda_{sell} \cdot E_{sell}
\end{equation}

其中每小时运维成本为：
\begin{equation}
C_{OM,h} = c_{ALKEL} \cdot P_{ALKEL}^{max} + c_{PEMEL} \cdot P_{PEMEL}^{max} + c_{NH3} \cdot P_{NH3}^{max} = 5003 \text{ 元/h}
\end{equation}

\subsubsection{求解方法}

采用PuLP建模框架调用CBC求解器进行精确求解。CBC 求解器基于 Branch-and-Cut 算法，通过线性松弛、变量分支、界限剪枝和割平面生成等步骤搜索整数可行解，最终得到全局最优调度方案。
该算法是在传统分支定界法的基础上结合割平面技术形成的混合整数规划求解方法。其核心思想是通过线性规划松弛获得界限，并利用分支定界搜索整数可行解；同时在搜索过程中加入割平面约束，切除不满足整数结构的松弛解，从而提高求解效率并保证全局最优性。对于每个产量水平$Q \in \{72, 63, 54, 45, 36\}$吨/日，分别构建MILP模型并求解最优开机时段组合。

\subsection{问题二(1)：典型场景各产量最优方案}

在典型风光场景下，对五种产量水平分别求解MILP，得到最优生产时段安排及对应的运行指标，结果如表\ref{tab:q2_typical}所示。

\begin{table}[H]
\centering
\caption{典型场景下各产量水平最优方案}\label{tab:q2_typical}
\begin{tabular}{cccccc}
\toprule
日产量(吨) & 开机时数 & 吨氨成本(元/吨) & $R_{self}$ & $R_{green}$ & $R_{grid}$ \\
\midrule
72 & 24h & 6219.14 & 86.5\% & 53.3\% & 6.7\% \\
63 & 21h & 5328.06 & 84.3\% & 59.7\% & 7.8\% \\
54 & 18h & 4591.27 & 80.2\% & 67.3\% & 9.9\% \\
45 & 15h & 3944.54 & 51.0\% & 66.7\% & 24.5\% \\
36 & 12h & 3191.67 & 10.5\% & 59.7\% & 44.8\% \\
\bottomrule
\end{tabular}
\end{table}

从吨氨成本角度看，产量越低则吨氨成本越低，36吨/日对应最低吨氨成本3191.67元/吨。这是因为低产量方案可以将开机时段集中在风光出力充足的白天，大幅减少购电需求。然而从绿电指标角度看，36吨/日方案的三项指标全部不满足要求（$R_{self}=10.5\%$远低于60\%，$R_{grid}=44.8\%$远超20\%），原因在于开机时段过短（仅12小时），集中在白天导致大量余电上网。

\begin{figure}[H]
\centering
\includegraphics[width=0.85\textwidth,height=0.38\textheight,keepaspectratio]{../figures/fig_q2_cost_vs_production.pdf}
\caption{不同日产氨量下的吨氨成本变化趋势}\label{fig:q2_cost_vs_production}
\end{figure}

吨氨成本随产量增加呈单调递增趋势（图\ref{fig:q2_cost_vs_production}），从36吨/日的3191.67元/吨上升至72吨/日的6219.14元/吨。这一趋势的经济学解释是：产量增加意味着更多时段需要开机，被迫在风光出力不足的夜间高电价时段运行，购电成本大幅增加。

\begin{figure}[H]
\centering
\includegraphics[width=0.9\textwidth,height=0.38\textheight,keepaspectratio]{../figures/fig_q2_schedule_gantt.pdf}
\caption{各产量水平下最优开停机时段安排}\label{fig:q2_schedule_gantt}
\end{figure}

最优开机时段安排呈现明显规律（图\ref{fig:q2_schedule_gantt}）：低产量方案优先选择白天光伏高峰时段开机，随着产量增加逐步向夜间扩展。绿电指标全满足的最优产量为54吨/日（$C_{ton}=4591.27$元/吨），此时$R_{self}=80.2\%>60\%$、$R_{green}=67.3\%>30\%$、$R_{grid}=9.9\%<20\%$，三项指标均满足政策要求。该产量对应18小时开机，设备利用率为75\%。

\subsection{问题二(2)：24种场景全年统计分析}

将6种风电出力水平与4种光伏出力水平组合，形成24种风光出力场景。对每种场景分别求解各产量水平的MILP，选取吨氨成本最低的方案作为该场景的最优运行策略。

\begin{figure}[H]
\centering
\includegraphics[width=0.85\textwidth,height=0.38\textheight,keepaspectratio]{../figures/fig_q2_scenarios_heatmap.pdf}
\caption{24种风光场景下吨氨成本热力图}\label{fig:q2_scenarios_heatmap}
\end{figure}

24种场景下的吨氨成本分布呈现显著差异（图\ref{fig:q2_scenarios_heatmap}）。高风电、高光伏场景（如W1P1）的吨氨成本最低（2055.78元/吨），因为风光出力充足可大幅减少购电；低风电、低光伏场景（如W6P4）的吨氨成本最高（7125.16元/吨），几乎全部依赖购电。风电出力水平对成本的影响大于光伏，这与风电装机容量（40 MW）相对于光伏（64 MW）的出力稳定性有关。

按照每种场景代表15天的假设，全年360天的绿电直连指标满足情况统计如表\ref{tab:q2_annual}所示。

\begin{table}[H]
\centering
\caption{问题二全年绿电直连指标统计}\label{tab:q2_annual}
\begin{tabular}{lcc}
\toprule
类别 & 场景数 & 对应天数 \\
\midrule
全满足（三项均达标） & 0种 & 0天 \\
部分满足（1--2项达标） & 21种 & 315天 \\
全不满足（三项均不达标） & 3种 & 45天 \\
\bottomrule
\end{tabular}
\end{table}

在离散调节模式下，全年没有任何场景能够同时满足三项绿电直连指标要求。这一结果说明，仅靠开停机调度无法有效解决风光出力与用电负荷的时序匹配问题，需要引入更灵活的功率调节手段。

\begin{figure}[H]
\centering
\includegraphics[width=0.9\textwidth,height=0.38\textheight,keepaspectratio]{../figures/fig_q2_annual_cost.pdf}
\caption{全年吨氨成本分布曲线}\label{fig:q2_annual_cost}
\end{figure}

全年吨氨成本分布曲线（图\ref{fig:q2_annual_cost}）显示，24种场景的吨氨成本从1370.40元/吨到7125.16元/吨不等，全年加权平均吨氨成本为4493.82元/吨，全年总产量为12960吨。成本分布呈现较大离散度，反映了风光出力随机性对园区经济性的显著影响。
